\section{AegisKV Design}
\label{sec:design}

AegisKV allocates finite KV storage to reduce workflow completion time
(WCT) while preserving normal cache reuse. Two observations connect
workflow-level priorities to the state needed for recovery.

\subsection{Observations}
\label{sec:design-insights}

\paragraph{1. Workflow dependencies determine recovery urgency.}
Figure~\ref{fig:motivation-recovery-exposure} shows similar reconstruction
costs with different effects on WCT: slack absorbs some branches' delays,
while a critical branch exposes its full pause. Thus request recovery
time alone is insufficient to rank protection priorities; the request's
position in the remaining workflow also matters.

\paragraph{2. Historical KV accesses concentrate on a subset of blocks.}
Figure~\ref{fig:sparse-releases} surveys long-context releases with block,
token or compressed-entry selection. Figure~\ref{fig:design-insights}
then examines native access traces: approximately 20\% of historical
blocks account for 87.64\%, 74.76\% and 56.54\% of accesses across the
three models. These results distinguish the complete historical state
from the subset frequently used by computation. The survey does not
establish a monotonic increase in sparsity, and retrospective concentration
does not establish future-query prediction accuracy.

\begin{figure}[t]
 \centering
 \includegraphics[width=\columnwidth]{figures/sparse-models.pdf}
 \caption{Selected sparse/hybrid model releases and reported context
 lengths. IDs: 1 MiniCPM4~\cite{minicpm4model};
 2 MiniCPM4.1~\cite{openbmb2025minicpm41};
 3 DeepSeek-V3.2-Exp~\cite{deepseek2025v32exp};
 4 DeepSeek-V3.2~\cite{deepseek2025v32model};
 5 GLM-5~\cite{zai2026glm5model};
 6 MiniCPM-SALA~\cite{openbmb2026sala};
 7 DeepSeek-V4 preview~\cite{deepseek2026v4model};
 8 MiniMax-M3~\cite{minimax2026m3};
 9 GLM-5.2~\cite{zai2026glm52};
 10 LongCat-Flash-Lite-Sparse~\cite{longcat2026sparse}.
 Dates are announcements, or first official repository dates for
 IDs 5, 9 and 10. Markers identify the sparse branch's selection unit,
 including SALA's block branch and V4's compressed sparse branch;
 they do not imply uniform attention across layers.}
 \label{fig:sparse-releases}
 \Description{Ten selected sparse or hybrid attention releases between
 June 2025 and July 2026, with reported contexts from 32K to 1M.
 Circles, squares and diamonds distinguish block, token and compressed
 entry selection. This is a literature survey, not a sparsity-rate curve.}
\end{figure}

\begin{figure}[t]
 \centering
 \includegraphics[width=\columnwidth]{figures/kv-access-concentration.pdf}
 \caption{Native KV access concentration for NOSA-8B~\cite{huang2026nosav2},
 NSA-8B (MiniCPM-4-8B-NSA)~\cite{fasa2026minicpm4nsa}, and
 MiniCPM4.1-8B~\cite{openbmb2025minicpm41}. We use the final 512 queries
 of matched 32,704-token prefixes from three PG19 texts. Each query
 counts a historical 64-token block once across heads. Rank the 503
 blocks within each layer/window; average access shares over all
 32 layers, then texts. Bands show text min--max; the diagonal denotes
 uniform access. Compressed/selector state is excluded. This is
 retrospective concentration, not prediction accuracy or recovery speedup.}
 \label{fig:design-insights}
 \Description{Cumulative historical block access curves for NOSA-8B,
 NSA-8B and MiniCPM4.1-8B, averaged over every layer and three texts.
 About 20 percent of blocks account for 87.64, 74.76 and 56.54 percent
 of accesses, respectively; all three curves exceed the uniform baseline.}
\end{figure}

Together, these observations guide which states to protect and which
parts to keep immediately accessible.

\subsection{System Overview}
\label{sec:system-overview}

AegisKV combines workflow priorities with a three-tier KV hierarchy
(Figure~\ref{fig:system-overview}). Shared DRAM retains complete checkpoints
for reuse and recovery; shared storage extends retention when DRAM is
under pressure. Peer HBM keeps selected layers and hot blocks on another
data-parallel (DP) replica to reduce recovery waiting. This division lets
the system retain long histories while spending scarce HBM on the state
whose faster access is most useful to workflow completion.

The Application supplies requests and workflow progress to the Controller.
\textbf{(1)} Workflow Tracker estimates how much delay each Agent can
absorb, while \textbf{(2)} Recovery-aware Router selects its execution and
takeover DPs. Using this urgency and location information, \textbf{(3)} KV
Master decides where to retain each checkpoint and which hot blocks merit
Peer space. Within the Engine, \textbf{(4)} Adapter adjusts the available
Peer capacity as serving demand changes, and \textbf{(5)} Connector moves
the selected KV. If a worker fails, Router redirects its requests to the
takeover DP, where Connector combines Peer hits with lower-tier reads to
resume execution.

\begin{figure}[t]
 \centering
 \includegraphics[width=\columnwidth]{figures/system-overview.pdf}
 \caption{AegisKV architecture. Callouts (1)--(5) match the text;
 DP1 mirrors DP0's Adapter/Connector. DP0 executes; DP1 hosts Peer
 replicas and takes over. Blue links carry KV, dark dashes control,
 red dashes retry; orange blocks are unsealed tails. Complete shared
 backing lies outside both DP failure boundaries.}
 \label{fig:system-overview}
 \Description{Compact Application, Controller and Engine layers with
 numbered components, symmetric DP Local/Peer HBM and shared backing.}
\end{figure}

\subsection{Workflow Tracker: Remaining Time and Slack}
\label{sec:workflow-protection}

To make workflow priorities available online, Tracker estimates remaining
execution before deriving slack. It maintains a residual DAG whose nodes
are Agent segments between exposed fork/join boundaries or completion.
A segment can span several LLM calls and tool waits. Child execution
appears on child paths, not again in the parent; unknown future forks
are excluded.

A GBDT predicts segment seconds remaining, $d$, from causal features: role/model,
current input length, elapsed time, completed-call count and previous
calls' lengths, timings and context growth. Missing history is masked;
training, calibration and evaluation split whole workflows. Labels end
at the segment boundary, and future call outcomes are never features.

Algorithm~\ref{alg:workflow-tracker} takes residual DAG $G$ and features
$X$. Forward/backward critical-path passes~\cite{kelley1959cpm} compute
earliest finishes $EF$ and latest finishes $LF$ relative to the predicted
workflow finish. Their difference is nominal slack. The guard $g$ is an
upper empirical quantile of positive slack overestimation on calibration
workflows; subtracting it gives advisory slack $s$, not a per-request
guarantee. Requests in one segment share its margin.

\par\medskip\noindent\begin{minipage}{\columnwidth}
\raggedright
 \refstepcounter{designalgorithm}\label{alg:workflow-tracker}
 \textbf{Algorithm \thedesignalgorithm: Remaining time and slack}
 \par\smallskip
 \begin{tabularx}{\columnwidth}{@{}r@{\quad}X@{}}
 \toprule
 1 & $(G,X)\gets\textsc{SnapshotProgress}()$ \\
 2 & $\nu\gets\textsc{NextVersion}()$ \\
 3 & \textbf{if not} \textsc{Supported}$(G,X)$ \textbf{then} \\
 4 & \quad \textsc{PublishInvalid}$(\nu)$; \textbf{return} \\
 5 & $d\gets\max(0,\textsc{GBDT}(X));\quad g\gets\textsc{Guard}()$ \\
 6 & $EF\gets\textsc{Forward}(G,d)$ \\
 7 & $LF\gets\textsc{Backward}(G,d,\max_v EF_v)$ \\
 8 & \textbf{for} $v\in V(G)$ \textbf{do} \\
 9 & \quad $s_v\gets\max(0,LF_v-EF_v-g_v)$ \\
 10 & \textsc{Publish}$(\nu,G,d,g,s)$ \\
 \bottomrule
 \end{tabularx}
\end{minipage}\par\medskip

\textsc{Supported} checks inputs, dependencies, calibration and freshness;
invalid estimates grant no positive slack credit. Progress events and
periodic refresh publish versioned durations $d$, guards and slack to
Master, invalidating older placement costs. The DAG passes cost
$O(|V|+|E|)$; prediction quality is evaluated in
Section~\ref{sec:eval-prediction}.

\subsection{Recovery-aware Router}
\label{sec:recovery-router}

Router establishes where normal and resumed execution run. Normal
dispatch balances load and compatible prefix locality.
Figure~\ref{fig:agent-prefix-reuse} reports 97.90\% potential same-agent
token coverage and 0.37\% additional cross-agent coverage. Agent affinity
is a preference: shared prefixes can still help another Agent's first call.

\begin{figure}[t]
 \centering
 \includegraphics[width=\columnwidth]{figures/agent-prefix-reuse.pdf}
 \caption{Potential prefix reuse over 98,827 requests, 2,090 agents and
 393 CC/Weka workflows. Coverage uses earlier completed requests with
 the same recorded model and workflow. Same/other-agent coverage
 overlaps; only additional cross-agent coverage extends same-agent
 reuse. These are token-prefix opportunities, not runtime cache hits.}
 \label{fig:agent-prefix-reuse}
 \Description{Same-agent prefixes dominate aggregate reuse; cross-agent
 history can also help first subagent calls.}
\end{figure}

At arrival, Router records a versioned takeover binding to a DP with a
compatible layout outside the source failure domain. Its corresponding
tensor-parallel (TP) shards host Peer copies. On confirmed failure, Router fences the old epoch and activates
that target subject to workspace and admission constraints. Rebinding
preserves a valid backing path and starts with zero assumed Peer coverage
on the new target. This binding supplies Master with a concrete location
for evaluating Peer access, transfer paths and target load.

\subsection{KV Master: Three-Tier Renewable Credit}
\label{sec:kv-master}

Using Tracker's workflow estimates and Router's target binding, Master
values backing placement, selects Peer layers/blocks, then applies those
choices through an online admission and eviction loop.

\subsubsection{Valuing complete backing and Peer acceleration}
\label{sec:placement-value}

Placement must account for both failure recovery and normal reuse.
The prefix opportunities in Figure~\ref{fig:agent-prefix-reuse} persist
across calls: repeatedly reconstructing a high-slack Agent's history
still consumes shared serving capacity. Slack therefore limits the
workflow benefit of faster recovery, but does not erase normal-reuse value.

\paragraph{Tier-conditional value.}
Each tier improves on the recovery path beneath it. A complete checkpoint
in storage avoids using an older checkpoint or reconstructing history.
DRAM shortens access, or provides retention when no storage copy exists;
Peer removes exposed read/H2D waiting.
All checkpoint dependencies remain in the chosen backing tier, including
blocks absent from Peer. Reclaiming DRAM therefore requires space for
the complete checkpoint in storage and revaluation of the changed Peer
miss path. Temporary copies share one retention benefit.
Figure~\ref{fig:retention-value} illustrates why ranking also requires
the amount of time saved and the bytes occupied.

\begin{figure}[t]
 \centering
 \includegraphics[width=\columnwidth]{figures/retention-value.pdf}
 \caption{Slack alone does not rank retention value. In this isolated-fault
 example, A uses two capacity units to reduce recovery from 4 to 2\,s;
 B uses one to reduce 18 to 8\,s. A at slack 0 saves 1\,s/unit, whereas
 B at slack 6 saves 10\,s/unit. Normal-use value is zero. These constructed
 profiles motivate cost/size comparison, not a policy-speedup claim.}
 \label{fig:retention-value}
 \Description{A higher-slack object can save more workflow time per byte.}
\end{figure}

\paragraph{Complete-plan objective.}
A common fault interrupts all requests in its source domain. At a join,
accelerating only one of two equally delayed branches need not improve
WCT. Master therefore compares complete placements, including victims
and moves. For current placement $P$ and candidate $P'$, it evaluates
\begin{equation}
 \Delta J=\sum_e p_e\sum_w[T_w(P',e)-T_w(P,e)].
 \label{eq:value}
\end{equation}
$T_w$ is predicted remaining workflow completion time. A fixed scenario
table covers normal continuation and source-domain faults at given offsets;
its probabilities sum to one. Both plans use the same queues, service
profiles, finite workspace and Connector order. Completed copies determine
the saved prefix and replay tail. Delays enter the residual DAG against
the same ordinary-execution reference, with joins taking the maximum.
Copy interference and temporary occupancy are included in this schedule.

The Tracker guard is added once per interrupted segment in both plans.
Invalid workflow timing falls back to summed normal/recovery request
completion delays without slack credit; invalid resource profiles pause
optional moves and use bounded recency order for safe mandatory relief.

\paragraph{Prefix credit for capacity competition.}
Credit screens candidates before joint evaluation. For a fixed history
and tail, Master adds groups from the tier's fallback in prefix order:
complete token extents across all layers, TP shards and auxiliary state
for backing, or layer/hot-block groups for Peer. Each step can retain
its previous restore path, keeping gains nonnegative. Credit combines
normal-use service-time savings and recovery delay saved beyond guarded
slack under the objective's horizon and scenarios. Reuse forecasts use
visible history; existing local hits earn no extra backing benefit.

Each increment is split by new object bytes. Distinct Agents contribute
separate shares; repeated references by one Agent do not duplicate its
share. Checkpoints can share unchanged blocks, charging capacity once per
physical copy. For example, restore times
12, 7 and 5\,s at successive frontiers with slack 6\,s yield failure-value
increments 5 and 1\,s before probability weighting, rather than assigning
all 6\,s to the tail. Increments telescope under the Agent's local model;
their sum across Agents screens competition but does not replace the
joint WCT objective.

\subsubsection{Choosing per-request layers and hot blocks}
\label{sec:sparse-placement}

Master next selects the Peer blocks that shorten recovery. A recovering
request reuses the $M$ tokens covered by its last completed checkpoint
and replays the $N$ acknowledged tokens generated since that checkpoint.
The saved prefix avoids historical reconstruction, but its KV still
requires lower-tier transfer. Computation on the replay tail can overlap
transfers for later layers. Peer is useful where those transfers would
otherwise leave the request waiting; its effective depth therefore
depends on the request's saved history, replay tail and recovery queues.

For request $r$, $L_r$ is the Peer prefix depth and $K_{r,l}$ the number
of token blocks in layer $l\le L_r$, including required TP objects.
Mandatory native dependencies come first; remaining blocks are ranked
by selection frequency over completed queries, counting each physical
block once per query across heads. Ties favor recency, then block ID.
A budget below the mandatory set is infeasible. Selection units, engine
blocks and kernel tiles need not coincide; the measured 64-token unit
is not a universal requirement. Complete backing supplies misses without
changing native selections (Figure~\ref{fig:sparse-recovery}).

\begin{figure}[t]
 \centering
 \includegraphics[width=\columnwidth]{figures/sparse-recovery.pdf}
 \caption{Complete backing, per-layer Peer budgets and demand misses.
 (a) The DRAM manifest covers B0--B3 in every layer; the tail is uncommitted.
 (b) On the bound target, $L=2$, with $K_1=2$ and $K_2=3$; layer 3 has
 no Peer copy. (c) Layer 1 needs B0/B2: B0 hits Peer and B2 is fetched
 from backing before computation. B3 was unused. The complete manifest
 also covers TP shards and selector state. If demoted as a whole,
 shared storage supplies the miss path before H2D.}
 \label{fig:sparse-recovery}
 \Description{Three complete historical layers, an uncommitted tail,
 different hot-block counts and an explicitly routed layer-1 miss.}
\end{figure}

Figure~\ref{fig:design-pipeline} isolates Peer's benefit: (c) and (d) share
complete DRAM backing and the same replay tail; Peer removes 24\,ms of
exposed waiting in this analytical instance. Its seven-layer depth is
profile-specific, motivating per-request choices rather than a fixed
depth for every request.

\begin{figure}[t]
 \centering
 \includegraphics[width=\columnwidth]{figures/recovery-pipeline.pdf}
 \caption{Analytical layered recovery with and without Peer.
 MiniCPM4-8B has 32 layers, 32K history and a
 64-token tail. Replay is calibrated at 2.870\,ms/layer; assumed effective
 H2D bandwidth is 6.25\,GB/s (5.369\,ms per full layer). Trace-derived
 mean preload and preload-plus-miss fractions are 0.6584/0.7748.
 (a) waits for bulk restoration; (b) overlaps full-layer restoration;
 (c) preloads hot blocks and fetches revealed misses; (d) adds 147.5\,MiB
 of Peer hot sets in the first seven layers. Numbers are tail completion
 and exposed wait. Cold copies follow replay; the ideal preemptible
 stream and homogeneous layer costs omit bounded-workspace and queue
 effects. This is not end-to-end fault-recovery measurement.}
 \label{fig:design-pipeline}
 \Description{The same required replay computation with progressively
 reduced exposed KV waits: 172, 83, 44 and 20 milliseconds.}
\end{figure}

Master enumerates a finite set of per-layer block-count or fraction
templates and depths from zero to model depth, plus the current placement.
This avoids a Cartesian product of all layer budgets. Bounded bundles
can jointly advance tied join branches; all evaluations share one work
budget. Each scenario uses current $M$ and its own predicted $N$,
including commit lag. Unsupported sparse profiles use full demand;
only computation that must still execute hides transfer.
Neither $L$ nor $K$ must increase with every token: a fixed budget can
change block identities, and $L=0$ is valid when no exposed wait is removed.
The online loop below decides when these candidates replace resident state.

\subsubsection{Online renewal, admission, and eviction}
\label{sec:placement-plan}

As calls arrive and end, the Agents interested in a prefix and their
benefits change. A one-time score misses these updates; renewing unchanged
prefixes at every checkpoint can instead preserve stale credit.
Master distinguishes actual reuse from value updates through renewable
credit and byte-proportional pressure aging, inspired by Landlord and
GreedyDual-Size~\cite{young2002landlord,cao1997greedydual}.

A physical replica $x$ has size $b_x$, each Agent's current grant $g_{a,x}$
and paid rent $d_{a,x}$. Each pool has its own cumulative rent $\rho_j$;
$L_r$ denotes only layer depth. Remaining credit and the eviction key are
\begin{equation}
 C_x=\sum_a[g_{a,x}-d_{a,x}]^+,
 \qquad H_x=\rho_j+C_x/b_x.
 \label{eq:rental-credit}
\end{equation}
Clocks and new paid amounts start at zero. Advancing rent by $\delta$
charges every committed replica $\min(C_x,b_x\delta)$, split proportionally
among its positive Agent shares and added to their paid amounts. Pinned
and leased copies also age; eligibility for removal is separate. Settlement
can be lazy. Rent represents capacity competition, not wall time.

An actual prefix/recovery use settles rent and renews only that Agent's
share by refreshing its grant and resetting its paid amount. Request/use
IDs prevent repeated renewal within one use. Value-only changes refresh
grants but preserve paid rent; checkpoint growth does not renew unchanged
prefixes. Expired demand removes its grant, not the copy. Figure~\ref{fig:renewable-credit}
illustrates the distinction.

\begin{figure}[t]
 \centering
 \includegraphics[width=\columnwidth]{figures/renewable-credit.pdf}
 \caption{Constructed credit accounting for one shared physical object
 of unit size. Initial Agent grants are 3 and 2. A rent charge of 2 leaves
 shares 1.8 and 1.2. Updating A's grant to 4 preserves its paid amount
 1.2, leaving 2.8; an actual new use by A then renews its share to 4.
 B's share remains 1.2 in both updates. Credit is a screening quantity,
 not measured WCT or duplicated physical capacity.}
 \label{fig:renewable-credit}
 \Description{Two Agents contribute to one physical prefix object.
 A four-row ledger separates pressure aging, value refresh and renewal.}
\end{figure}

\paragraph{Legal competition.}
Victims are legal moves: complete-checkpoint demotion or fallback,
or valid Peer hot-set reduction. Bytes remain occupied while another
checkpoint references them; bounded owner bundles can release shared objects.
Ordinary competition respects running requests' logical Peer pins.
Request end/tool handoff unpins; self replacement and required capacity
relief may override logical pins, never physical read/copy leases.

For each move, \textsc{RentPlans} refreshes up to $q\ge2$ low-key victims
and $q$ rotating-scan alternatives, exploring victim chains within a common
work budget. Releasing objects $Y$ raises the trial clock to at least
$\max_{x\in Y}H_x$. A group of new bytes must retain positive credit:
$G_{\rm in}>b_{\rm in}(\rho'_j-\rho_j)$, where $G_{\rm in}$ and
$b_{\rm in}$ are the new bytes' grant and size, and $\rho'_j$ is the trial
clock. The trial charge is split by incoming grants; existing replicas neither consume new capacity nor gain
new credit. Each pool applies its own gate. Plans include transient space
and recursively secure storage for DRAM demotion; Peer release uses its
existing backing without redundant D2H.

Prefill offers newly sealed state; decode triggers review every fixed,
block-aligned $Z$ tokens. Algorithm~\ref{alg:placement} also handles use,
value, request-end, budget, failure and transfer events $e$.
Each event uses a state snapshot $S$; $h$ marks required capacity relief.
Candidates include complete checkpoint extensions (DRAM first, then
storage), promotion/demotion, last-copy removal with a valid fallback,
and Peer replacements from Section~\ref{sec:sparse-placement}.
If neither backing tier admits an extension, $M$ stays unchanged;
copy lag can make $N$ exceed $Z$.

\par\medskip\noindent\begin{minipage}{\columnwidth}
\raggedright
 \refstepcounter{designalgorithm}\label{alg:placement}
 \textbf{Algorithm \thedesignalgorithm: Three-tier rental event loop}
 \par\smallskip
 \begin{tabularx}{\columnwidth}{@{}r@{\quad}X@{}}
 \toprule
 1 & \textsc{Reconcile}$(e)$; \textsc{SettleAffected}$(e)$ \\
 2 & \textbf{if} \textsc{NewUse}$(e)$ \textbf{then} \textsc{RenewAgent}$(e)$ \\
 3 & \textsc{MarkChangedValues}$(e)$; $S\gets\textsc{Snapshot}()$ \\
 4 & $h\gets\textsc{ShrinkRequired}(S)$; $(P^*,d^*)\gets(\varnothing,\infty)$ \\
 5 & \textbf{for} $m\in\textsc{Candidates}(S,h)$ \textbf{while} \textsc{WorkLeft}() \\
 6 & \quad \textbf{for} $P\in\textsc{RentPlans}(S,m,h)$ \\
 7 & \qquad $d\gets\Delta J(P;S)$ \\
 8 & \qquad \textbf{if} $d<d^*$ \textbf{then} $(P^*,d^*)\gets(P,d)$ \\
 9 & \textbf{if} $P^*=\varnothing$ \textbf{then} \textsc{Defer}$(e)$; \textbf{return} \\
 10 & \textbf{if} $\neg h\land d^*\ge-\epsilon$ \textbf{then return} \\
 11 & \textbf{if not} \textsc{ReserveIfCurrent}$(S,P^*)$ \textbf{then} \\
 12 & \quad \textsc{Requeue}$(e)$; \textbf{return} \\
 13 & \textsc{Submit}$(P^*)$ \quad\textit{completion/failure emits a new event} \\
 \bottomrule
 \end{tabularx}
\end{minipage}\par\medskip

\textsc{Reconcile} applies verified commits, unpins and failures;
value refresh preserves paid rent. Optional
plans must pass rent screening and $\Delta J<-\epsilon$. Required relief
bypasses positive-credit admission and chooses the least costly examined
legal plan, still charging released objects. If no safe plan is found,
admission pauses until a later event retries the newest pending frontier.
Periodic budgeted refresh revisits stale keys. Tool waits retain Agent
interest; terminal progress removes failure demand, and expired inactive
shares are pruned.

Copies may outlive the decision that scheduled them. Before acting,
\textsc{ReserveIfCurrent} therefore rechecks versions, leases, bytes and pool clocks.
Rejected previews change no committed rent or residency. A prepared
handoff advances rent only after its fallback is published; new replicas
install their residual trial credit on commit. Failure before handoff
discards trial rent, whereas failure after handoff cannot undo released
victims or reduce the clock. Connector completion/failure feeds the next
event, closing the loop without treating a pending copy as protection.
This bounded adaptation offers no optimality or inherited competitive guarantee.

\subsection{Adapter: Guarded AIMD for Peer Capacity}
\label{sec:adapter}

Master selects content within a quota; Adapter adjusts that quota as
foreground inference demand changes. It uses pressure-driven multiplicative
decrease and stable-window additive increase~\cite{chiu1989aimd}.
The total Peer authorization has a physical ceiling
\begin{equation}
 C=\max\{0,\min(B_{\mathrm{cap}},
 B_{\mathrm{pool}}-B_{\mathrm{serve}}-B_{\mathrm{reserve}})\}.
 \label{eq:peer-ceiling}
\end{equation}
The pool excludes weights/fixed allocations; serving high watermark
excludes Peer, and reserve covers growth/workspace. Existing Peer bytes
are not subtracted twice. $B_{\rm cap}$ can be a configured pool fraction.

The current quota $B$ and stable-window count $z$ start at zero.
Pressure (allocation stalls, reserve violations or
memory-related admission waits) multiplies $B$ by $0<\gamma<1$. After $m$
stable windows below a lower wait threshold, deferred positive-value Peer
demand, assessed before quota filtering, permits an additive byte step
$\Delta$. Intermediate windows reset
$z$; missing telemetry disables growth. The ceiling clips every update;
emergencies may decrease within a window, but growth requires completed
stable windows. Parameters and window length are configurable.

\par\medskip\noindent\begin{minipage}{\columnwidth}
\raggedright
 \refstepcounter{designalgorithm}\label{alg:adapter}
 \textbf{Algorithm \thedesignalgorithm: Windowed Peer quota}
 \par\smallskip
 \begin{tabularx}{\columnwidth}{@{}r@{\quad}X@{}}
 \toprule
 1 & $C\gets\textsc{SafeCeiling}();\quad B'\gets B$ \\
 2 & \textbf{if} \textsc{Pressure}() \textbf{then} $(B',z)\gets(\gamma B,0)$ \\
 3 & \textbf{else if} \textsc{Stable}() \textbf{then} \\
 4 & \quad $z\gets z+1$ \\
 5 & \quad \textbf{if} $z\ge m$ \textbf{then} \\
 6 & \qquad \textbf{if} \textsc{PeerDemand}() \textbf{then} $B'\gets B+\Delta$ \\
 7 & \qquad $z\gets0$ \\
 8 & \textbf{else} $z\gets0$ \\
 9 & $B\gets\textsc{RoundDownPage}(\min(C,B'))$ \\
 10 & \textsc{PublishBudget}$(B)$; \textsc{EmitCapacityEvent}() \\
 \bottomrule
 \end{tabularx}
\end{minipage}\par\medskip

Adapter sets the total budget; Master chooses $L/K$ at fixed $Z$ and
other placement events. If occupancy plus reservations exceeds a reduced
quota, Master pauses Peer admission and selects replicas to release.
Serving may wait for ongoing readers and transfers to finish before
reusing their memory.

\subsection{Connector: Committed Copies and Layered Recovery}
\label{sec:connector}

Connector executes placement changes alongside inference, sharing transfer
capacity among reuse, recovery and checkpoints through bounded queues.
Recovery misses precede speculative reads. A final snapshot at request
completion or tool handoff captures short rounds that never reach $Z$;
accepted copies retain their source KV until transfer completes.

While copying proceeds, the previous checkpoint remains usable.
Connector publishes the new checkpoint's block list, or manifest, only
when its required state has arrived and matches the token log. Moving
to storage additionally requires persistence before reads are redirected
and DRAM is reclaimed; existing readers finish before their memory is
reused. Failed or partial copies leave the old recovery point in place.

Accepted inputs and generated tokens enter an ordered log before output
acknowledgement, enabling replay beyond the checkpoint. The log, metadata
and backing remain outside the failed worker's domain. Last-copy eviction
first installs an older checkpoint or input-and-log replay path and
withdraws dependent Peer coverage, preserving a way to resume.

On the takeover DP, Peer hits are used directly; misses come from DRAM
or storage followed by H2D. Reads overlap tail replay, with each operation
waiting for its actual KV, selector and TP dependencies. Prediction misses
therefore cause extra fetches. Replaying acknowledged tokens beyond $M$
restores their KV without re-emitting outputs, after which decode resumes.
Adopted Peer blocks transfer to execution accounting without another allocation.
